% Gerado por ferramentas/gerar.py; não edite à mão.
\section*{Exercícios}
\addcontentsline{toc}{section}{Exercícios}
Os pontos antes de cada número indicam a dificuldade: \nivel{1} fácil, \nivel{2} médio, \nivel{3} difícil e \nivel{4} desafio. Os exercícios marcados com \estrela{} têm solução completa, com animação do gráfico, no repositório do livro: \url{https://github.com/AleffAzevedo/calculo-1-pt/blob/main/exercicios/cap7/README.md}. As respostas estão no fim do capítulo.

\begin{itemize}[leftmargin=5.2em,labelsep=0.6em,itemsep=4pt]
\item[\exnum{1}{7.1}] Calcule $\displaystyle\int x\cos x\,dx$ e $\displaystyle\int xe^{-x}\,dx$.
\item[\exnum{1}{7.2}] Calcule $\displaystyle\int x\ln x\,dx$.
\item[\exnum{1}{7.3}] Calcule $\displaystyle\int\tan^{-1}x\,dx$ (use $dv=dx$).
\item[\exnum{1}{7.4}] Calcule $\displaystyle\int_0^{\pi/2}x\sin x\,dx$.
\item[\exnum{1}{7.5}] Calcule $\displaystyle\int\sin^3x\,dx$.
\item[\exnum{1}{7.6}] Calcule $\displaystyle\int\cos^2x\,dx$ e $\displaystyle\int_0^\pi\sin^2x\,dx$.
\item[\exnum{1}{7.7}] Calcule $\displaystyle\int\sin^2x\cos^3x\,dx$.
\item[\exnum{1}{7.8}] Calcule $\displaystyle\int\tan^2x\,dx$ e $\displaystyle\int\sec^4x\,dx$.
\item[\exnum{1}{7.9}] Calcule $\displaystyle\int\sin3x\cos2x\,dx$ usando produto em soma.
\item[\exnum{1}{7.10}] Calcule $\displaystyle\int_0^2\sqrt{4-x^2}\,dx$ com $x=2\sin\theta$ e confira com a área de um quarto de círculo.
\item[\exnum{1}{7.11}] Calcule $\displaystyle\int\frac{dx}{(x^2+9)^{3/2}}$ com $x=3\tan\theta$.
\item[\exnum{1}{7.12}] Calcule $\displaystyle\int\frac{dx}{(x-1)(x+2)}$.
\item[\exnum{1}{7.13}] Calcule $\displaystyle\int\frac{3x+1}{x^2-1}\,dx$.
\item[\exnum{1}{7.14}] Calcule $\displaystyle\int\frac{x^3}{x-1}\,dx$ (divida primeiro).
\item[\exnum{1}{7.15}] Calcule $T_4$ e $S_4$ para $\int_0^2x^3\,dx$ e compare com o valor exato. Por que $S_4$ acerta?
\item[\exnum{1}{7.16}] Decida se convergem e calcule quando for o caso: $\int_1^\infty\frac{dx}{x^3}$ e $\int_1^\infty\frac{dx}{\sqrt x}$.
\item[\exnum{1}{7.17}] Calcule $\int_0^\infty e^{-2x}\,dx$ e $\int_0^1\ln x\,dx$.
\item[\exnum{1}{7.18}] Calcule $\int_0^4\frac{dx}{\sqrt{4-x}}$ e decida se $\int_0^2\frac{dx}{(x-1)^2}$ converge.
\item[\exnum{1}{7.19}] Calcule $\displaystyle\int x^2\sin x\,dx$ e confira derivando a resposta.
\item[\exnum{1}{7.20}] Calcule $\displaystyle\int e^x\cos x\,dx$ (integral cíclica).
\item[\exnum{2}{7.21}] \estrela\ \textbf{Partes com logaritmo (e a figura do retângulo).} Calcule $\int_1^e\ln x\,dx$ e $\int_1^ex\ln x\,dx$ por partes. Interprete o primeiro resultado como
uma diferença de áreas num retângulo.
\item[\exnum{2}{7.22}] \estrela\ \textbf{Uma integral que volta a si mesma.} Calcule $\int e^x\sin x\,dx$ e depois $\int_0^\pi e^x\sin x\,dx$.
\item[\exnum{2}{7.23}] \estrela\ \textbf{A integral da secante ao cubo.} Mostre que $\int\sec^3x\,dx=\frac12\big(\sec x\tan x+\ln|\sec x+\tan x|\big)+C$ e calcule
$\int_0^{\pi/4}\sec^3x\,dx$.
\item[\exnum{2}{7.24}] \estrela\ \textbf{A área da elipse.} Use substituição trigonométrica para mostrar que a elipse $\frac{x^2}{a^2}+\frac{y^2}{b^2}\le1$ tem área $\pi ab$.
\item[\exnum{2}{7.25}] \estrela\ \textbf{Frações parciais com fator repetido.} Calcule $\displaystyle\int\frac{5x^2+3x-2}{x^3+2x^2}\,dx$ e depois a integral de $1$ a $2$.
\item[\exnum{2}{7.26}] \estrela\ \textbf{Trapézio e Simpson para $e^{-x^2}$.} A função $e^{-x^2}$ não tem antiderivada elementar. Calcule $T_4$ e $S_4$ para $\int_0^1e^{-x^2}dx$ e
use os limitantes de erro (com $|f''|\le2$ e $|f^{(4)}|\le12$ em $[0,1]$) para dizer quantas casas são confiáveis.
\item[\exnum{2}{7.27}] Calcule $\displaystyle\int x^2\cos x\,dx$.
\item[\exnum{2}{7.28}] Calcule $\displaystyle\int x^3e^{x^2}\,dx$ (substituição e depois partes).
\item[\exnum{2}{7.29}] Calcule $\displaystyle\int\sin^{-1}x\,dx$ e $\displaystyle\int(\ln x)^2\,dx$.
\item[\exnum{2}{7.30}] Calcule $\displaystyle\int\sin^4x\,dx$.
\item[\exnum{2}{7.31}] Calcule $\displaystyle\int\tan^3x\sec^3x\,dx$.
\item[\exnum{2}{7.32}] Calcule $\displaystyle\int\frac{\sqrt{x^2-9}}{x}\,dx$ ($x>3$).
\item[\exnum{2}{7.33}] Calcule $\displaystyle\int\frac{x^2}{\sqrt{9-x^2}}\,dx$.
\item[\exnum{2}{7.34}] Calcule $\displaystyle\int\frac{x^2+1}{x(x-1)^2}\,dx$.
\item[\exnum{2}{7.35}] Calcule $\displaystyle\int\frac{4}{x^3+4x}\,dx$.
\item[\exnum{2}{7.36}] Calcule $\displaystyle\int\frac{dx}{e^x+1}$ (faça $u=e^x$ e use frações parciais).
\item[\exnum{2}{7.37}] Calcule $S_4$ para $\int_0^\pi\sin x\,dx$. Qual o menor $n$ que garante, pelo teorema, erro de Simpson menor que $10^{-6}$?
\item[\exnum{2}{7.38}] Calcule $\int_0^\infty xe^{-x}\,dx$ e $\int_{-\infty}^\infty\frac{dx}{x^2+4x+5}$.
\item[\exnum{2}{7.39}] Calcule $\int_0^3\frac{dx}{(x-1)^{2/3}}$ (cuidado com $x=1$).
\item[\exnum{2}{7.40}] Use comparação para decidir: $\int_1^\infty\frac{dx}{x^2+\sqrt x}$ e $\int_2^\infty\frac{dx}{\ln x}$.
\item[\exnum{3}{7.41}] \estrela\ \textbf{O comprimento da parábola.} Calcule o comprimento de arco de $y=x^2$ para $0\le x\le1$.
\item[\exnum{3}{7.42}] \estrela\ \textbf{A trombeta de Gabriel.} Gire $y=\frac1x$, $x\ge1$, em torno do eixo $x$. Mostre que o sólido tem volume finito, mas a
superfície tem área infinita.
\item[\exnum{3}{7.43}] \estrela\ \textbf{A função gama.} Seja $\Gamma(x)=\int_0^\infty t^{x-1}e^{-t}\,dt$ para $x>0$. Mostre que $\Gamma(1)=1$, que
$\Gamma(x+1)=x\,\Gamma(x)$ e conclua que $\Gamma(n+1)=n!$.
\item[\exnum{3}{7.44}] \estrela\ \textbf{A substituição universal $t=\tan(x/2)$.} Com $t=\tan\frac x2$, mostre que $\cos x=\frac{1-t^2}{1+t^2}$ e $dx=\frac{2\,dt}{1+t^2}$, e calcule
$\displaystyle\int_0^{\pi/2}\frac{dx}{2+\cos x}$.
\item[\exnum{3}{7.45}] \estrela\ \textbf{Crescimento logístico.} Uma população segue $\dfrac{dP}{dt}=0{,}5\,P\Big(1-\dfrac P{1000}\Big)$, com $P(0)=100$. Separe as variáveis,
use frações parciais e encontre $P(t)$. Quando a população chega a $500$? Quando cresce mais rápido?
\item[\exnum{3}{7.46}] Calcule $\displaystyle\int e^{2x}\sin3x\,dx$.
\item[\exnum{3}{7.47}] Calcule $\displaystyle\int\cos\sqrt x\,dx$.
\item[\exnum{3}{7.48}] Deduza a fórmula de redução $\int\sin^nx\,dx=-\frac1n\sin^{n-1}x\cos x+\frac{n-1}n\int\sin^{n-2}x\,dx$ e calcule $\int_0^{\pi/2}\sin^6x\,dx$.
\item[\exnum{3}{7.49}] Calcule $\displaystyle\int x\tan^{-1}x\,dx$.
\item[\exnum{3}{7.50}] Calcule $\displaystyle\int\frac{dx}{(1+x^2)^2}$ e $\displaystyle\int_0^1\frac{dx}{(1+x^2)^2}$.
\item[\exnum{3}{7.51}] Calcule $\displaystyle\int\frac{e^x}{e^{2x}+3e^x+2}\,dx$.
\item[\exnum{3}{7.52}] Calcule $\displaystyle\int\frac{x^3+x+1}{(x^2+1)^2}\,dx$.
\item[\exnum{3}{7.53}] Calcule $\displaystyle\int\frac{dx}{x^3+1}$.
\item[\exnum{3}{7.54}] Calcule $\displaystyle\int\frac{\sqrt x}{1+x}\,dx$ (racionalize com $u=\sqrt x$).
\item[\exnum{3}{7.55}] Quantos subintervalos garantem, pela regra do trapézio, erro menor que $10^{-4}$ em $\int_1^2\frac{dx}x$?
\item[\exnum{3}{7.56}] Calcule $\int_0^\infty e^{-x}\sin x\,dx$.
\item[\exnum{3}{7.57}] Calcule $\int_1^\infty\frac{\ln x}{x^2}\,dx$.
\item[\exnum{3}{7.58}] Para quais $p$ a integral $\int_0^1\frac{dx}{x^p}$ converge, e para quanto? E $\int_0^\infty\frac{dx}{x^p}$?
\item[\exnum{3}{7.59}] Calcule $\displaystyle\int_0^\infty\frac{dx}{(1+x)\sqrt x}$ (é imprópria nas duas pontas).
\item[\exnum{3}{7.60}] A transformada de Laplace de $f$ é $F(s)=\int_0^\infty e^{-st}f(t)\,dt$. Calcule-a para $f(t)=1$ e $f(t)=t$ ($s>0$).
\item[\exnum{4}{7.61}] Prove que $\displaystyle\int_0^1\frac{x^4(1-x)^4}{1+x^2}\,dx=\frac{22}7-\pi$ e conclua que $\pi<\frac{22}7$.
\item[\exnum{4}{7.62}] Calcule $\displaystyle\int_0^{\pi/2}\ln(\sin x)\,dx$. (Dica: compare com $\int_0^{\pi/2}\ln(\cos x)\,dx$ e use $\sin2x=2\sin x\cos x$.)
\item[\exnum{4}{7.63}] Prove que $\displaystyle\int_1^\infty\frac{\sin x}x\,dx$ converge (partes com $u=\frac1x$), mas $\int_1^\infty\frac{|\sin x|}x\,dx$ diverge.
\item[\exnum{4}{7.64}] Sabendo que $\int_0^\infty e^{-x^2}dx=\frac{\sqrt\pi}2$, mostre que $\Gamma(\tfrac12)=\sqrt\pi$ e calcule $\Gamma(\tfrac52)$.
\item[\exnum{4}{7.65}] Sejam $I_n=\int_0^{\pi/2}\sin^nx\,dx$. Mostre que $I_{2n+1}\le I_{2n}\le I_{2n-1}$ e use a fórmula de redução para provar o produto de Wallis $\frac\pi2=\lim_{n\to\infty}\frac{2\cdot2\cdot4\cdot4\cdots2n\cdot2n}{1\cdot3\cdot3\cdot5\cdots(2n-1)(2n+1)}$.
\end{itemize}
\subsection*{Respostas}
{\small\begin{multicols}{2}\raggedright
\begin{itemize}[leftmargin=2.7em,labelsep=0.5em,itemsep=1pt]
\item[\textbf{7.1}] $x\sin x+\cos x+C$; $-(x+1)e^{-x}+C$.
\item[\textbf{7.2}] $\tfrac{x^2}2\ln x-\tfrac{x^2}4+C$.
\item[\textbf{7.3}] $x\tan^{-1}x-\tfrac12\ln(1+x^2)+C$.
\item[\textbf{7.4}] $1$.
\item[\textbf{7.5}] $-\cos x+\tfrac13\cos^3x+C$.
\item[\textbf{7.6}] $\tfrac x2+\tfrac{\sin2x}4+C$; $\tfrac\pi2$.
\item[\textbf{7.7}] $\tfrac13\sin^3x-\tfrac15\sin^5x+C$.
\item[\textbf{7.8}] $\tan x-x+C$; $\tan x+\tfrac13\tan^3x+C$.
\item[\textbf{7.9}] $-\tfrac1{10}\cos5x-\tfrac12\cos x+C$.
\item[\textbf{7.10}] $\pi$.
\item[\textbf{7.11}] $\tfrac{x}{9\sqrt{x^2+9}}+C$.
\item[\textbf{7.12}] $\tfrac13\ln\Big|\tfrac{x-1}{x+2}\Big|+C$.
\item[\textbf{7.13}] $2\ln|x-1|+\ln|x+1|+C$.
\item[\textbf{7.14}] $\tfrac{x^3}3+\tfrac{x^2}2+x+\ln|x-1|+C$.
\item[\textbf{7.15}] $T_4=4{,}25$; $S_4=4$ (exato): Simpson é exata para cúbicas.
\item[\textbf{7.16}] $\tfrac12$; diverge.
\item[\textbf{7.17}] $\tfrac12$; $-1$.
\item[\textbf{7.18}] $4$; diverge (o integrando explode em $x=1$).
\item[\textbf{7.19}] $-x^2\cos x+2x\sin x+2\cos x+C$.
\item[\textbf{7.20}] $\tfrac12e^x(\sin x+\cos x)+C$.
\item[\textbf{7.21}] $1$; $\tfrac{e^2+1}4\approx2{,}097$.
\item[\textbf{7.22}] $\tfrac12e^x(\sin x-\cos x)+C$; $\tfrac{e^\pi+1}2\approx12{,}070$.
\item[\textbf{7.23}] $\tfrac12\big(\sqrt2+\ln(1+\sqrt2)\big)\approx1{,}1478$.
\item[\textbf{7.24}] $\pi ab$.
\item[\textbf{7.25}] $2\ln|x|+\frac1x+3\ln|x+2|+C$; $2\ln2-\tfrac12+3\ln\tfrac43\approx1{,}749$.
\item[\textbf{7.26}] $T_4\approx0{,}742984$; $S_4\approx0{,}746855$; valor $\approx0{,}746824$.
\item[\textbf{7.27}] $x^2\sin x+2x\cos x-2\sin x+C$.
\item[\textbf{7.28}] $\tfrac12e^{x^2}(x^2-1)+C$.
\item[\textbf{7.29}] $x\sin^{-1}x+\sqrt{1-x^2}+C$; $x(\ln x)^2-2x\ln x+2x+C$.
\item[\textbf{7.30}] $\tfrac{3x}8-\tfrac{\sin2x}4+\tfrac{\sin4x}{32}+C$.
\item[\textbf{7.31}] $\tfrac15\sec^5x-\tfrac13\sec^3x+C$.
\item[\textbf{7.32}] $\sqrt{x^2-9}-3\sec^{-1}\tfrac x3+C$.
\item[\textbf{7.33}] $\tfrac92\sin^{-1}\tfrac x3-\tfrac x2\sqrt{9-x^2}+C$.
\item[\textbf{7.34}] $\ln|x|-\tfrac2{x-1}+C$.
\item[\textbf{7.35}] $\ln|x|-\tfrac12\ln(x^2+4)+C$.
\item[\textbf{7.36}] $x-\ln(1+e^x)+C$.
\item[\textbf{7.37}] $S_4=\tfrac\pi{12}(4\sqrt2+2)\approx2{,}00456$; $n=38$.
\item[\textbf{7.38}] $1$; $\pi$.
\item[\textbf{7.39}] $3+3\sqrt[3]2\approx6{,}780$.
\item[\textbf{7.40}] Converge ($\le\tfrac1{x^2}$); diverge ($\tfrac1{\ln x}\ge\tfrac1x$).
\item[\textbf{7.41}] $\tfrac14\big(2\sqrt5+\ln(2+\sqrt5)\big)\approx1{,}4789$.
\item[\textbf{7.42}] $V=\pi$; $S=\infty$.
\item[\textbf{7.43}] $\Gamma(n+1)=n!$.
\item[\textbf{7.44}] $\tfrac{\pi}{3\sqrt3}\approx0{,}6046$.
\item[\textbf{7.45}] $P(t)=\dfrac{1000}{1+9e^{-t/2}}$; $t=2\ln9\approx4{,}39$; crescimento máximo quando $P=500$.
\item[\textbf{7.46}] $\tfrac1{13}e^{2x}(2\sin3x-3\cos3x)+C$.
\item[\textbf{7.47}] $2\big(\sqrt x\sin\sqrt x+\cos\sqrt x\big)+C$.
\item[\textbf{7.48}] $\tfrac56\cdot\tfrac34\cdot\tfrac12\cdot\tfrac\pi2=\tfrac{5\pi}{32}$.
\item[\textbf{7.49}] $\tfrac12\big[(x^2+1)\tan^{-1}x-x\big]+C$.
\item[\textbf{7.50}] $\tfrac12\big(\tan^{-1}x+\tfrac x{1+x^2}\big)+C$; $\tfrac\pi8+\tfrac14$.
\item[\textbf{7.51}] $\ln\tfrac{e^x+1}{e^x+2}+C$.
\item[\textbf{7.52}] $\tfrac12\ln(x^2+1)+\tfrac12\big(\tan^{-1}x+\tfrac x{1+x^2}\big)+C$.
\item[\textbf{7.53}] $\tfrac13\ln|x+1|-\tfrac16\ln(x^2-x+1)+\tfrac1{\sqrt3}\tan^{-1}\tfrac{2x-1}{\sqrt3}+C$.
\item[\textbf{7.54}] $2\sqrt x-2\tan^{-1}\sqrt x+C$.
\item[\textbf{7.55}] $\tfrac{2}{12n^2}<10^{-4}$: $n\ge41$.
\item[\textbf{7.56}] $\tfrac12$.
\item[\textbf{7.57}] $1$.
\item[\textbf{7.58}] $p<1$, vale $\tfrac1{1-p}$; a segunda diverge para todo $p$.
\item[\textbf{7.59}] $\pi$ (com $u=\sqrt x$).
\item[\textbf{7.60}] $\tfrac1s$; $\tfrac1{s^2}$.
\item[\textbf{7.61}] Divida: o quociente é $x^6-4x^5+5x^4-4x^2+4$ e o resto dá $-4\tan^{-1}x$; o integrando é positivo.
\item[\textbf{7.62}] $-\tfrac\pi2\ln2$.
\item[\textbf{7.63}] Partes: $\cos1-\int_1^\infty\tfrac{\cos x}{x^2}dx$, absolutamente convergente; e $|\sin x|\ge\sin^2x=\tfrac{1-\cos2x}2$.
\item[\textbf{7.64}] $t=x^2$; $\Gamma(\tfrac52)=\tfrac32\cdot\tfrac12\sqrt\pi=\tfrac{3\sqrt\pi}4$.
\item[\textbf{7.65}] $\tfrac{I_{2n}}{I_{2n+1}}\to1$ pelo sanduíche, e essa razão é $\tfrac\pi2$ dividido pelo produto parcial.
\end{itemize}\end{multicols}}
